\documentclass[11pt]{article}
\usepackage[a4paper, margin=1in]{geometry}
\usepackage[T1]{fontenc}
\usepackage[utf8]{inputenc}
\usepackage{amsmath, amssymb}
\usepackage{microtype}
\usepackage{parskip}
\usepackage{hyperref}
\hypersetup{colorlinks=true, urlcolor=blue}
\pagestyle{empty}

\begin{document}

\noindent
Claes Johnson, Professor emeritus of Applied Mathematics \\
KTH Royal Institute of Technology, Stockholm \quad\textbar\quad \texttt{claesjohnson@gmail.com}

\vspace{1em}

\noindent
To the Editors, \emph{Studies in History and Philosophy of Science} \hfill \today

\vspace{1em}

\noindent Dear Editors,

\medskip

I submit for your consideration the article \emph{Quantum Mechanics without Ontology? Configuration Space, the
Arena Problem, and the Road Not Taken in 1926}, for publication in \emph{Studies in History and Philosophy of
Science}. The paper sits squarely at the intersection of the history and the philosophy of modern physics that
the journal exists to serve: it combines a reading of the 1926 turn to configuration space with an argument
about the metaphysics of the quantum state.

Its contribution is threefold. \emph{First}, it sharpens the common claim that quantum mechanics ``abandons
ontology.'' Taking the Pusey--Barrett--Rudolph theorem seriously, the paper grants that the quantum state is
plausibly ontic, and relocates the difficulty from the \emph{existence} of an ontology to its \emph{arena}: the
state that PBR renders real is a field on $3N$-dimensional configuration space, not on the three-dimensional
space of observers and instruments. The measurement problem, collapse, the Heisenberg cut, and the nonlocality
puzzles are then presented as symptoms of this \emph{arena problem} rather than of a bare absence of the real.
\emph{Second}, it situates the resulting options explicitly within the contemporary debate between wave-function
realism (Albert, Ney, Wallace) and primitive-ontology approaches (Allori, D\"urr--Goldstein--Zangh\`i, Maudlin,
Bell), and it revisits the historical moment --- with reference to Schr\"odinger's correspondence with Lorentz
and to Mehra and Rechenberg --- at which the configuration-space arena was chosen over the objections of
Schr\"odinger and Einstein, arguing that the choice was contingent rather than forced. \emph{Third}, it offers a
concrete three-dimensional charge-density formulation as an under-explored member of the primitive-ontology
family --- an existence proof that the road not taken in 1926 remains open --- and draws a practical corollary
for quantum computing, whose claimed advantage turns on the physical reality of the exponentially large
configuration-space amplitude.

The paper is, I believe, of interest to historians and philosophers of physics alike: it uses a precise recent
result (PBR) and a live metaphysical debate (wave-function realism versus primitive ontology) to give the
century-old interpretive impasse a clearer diagnosis, and it grounds the philosophical claim in a datable
historical episode. The manuscript is original, has not been published elsewhere, and is not under consideration
by another journal.

\medskip

\noindent In the interest of full transparency: the manuscript was prepared in cooperation with an AI assistant
(Claude, Anthropic), which contributed to drafting and to surveying the secondary literature; the argument,
claims, and responsibility for the content are the author's. I am glad to describe this collaboration further if
the journal wishes.

\medskip

\noindent Thank you for your consideration.

\vspace{1em}

\noindent Sincerely, \\
Claes Johnson

\end{document}
